\chapter{Liver biopsy}
\index{Liver biopsy}

\section*{Definition and Overview}
A liver biopsy is a medical procedure in which a small sample of liver tissue is removed using a special needle so it can be examined under a microscope. This diagnostic tool is used to identify liver diseases, assess the extent of liver damage (such as fibrosis or cirrhosis), and monitor the effectiveness of treatments. The procedure can be performed percutaneously (through the skin), transjugularly (through the neck vein), or laparoscopically during surgery.

\section*{Epidemiology}
Liver biopsies are performed in patients of all ages who present with unexplained abnormal liver function tests, chronic liver disease, or liver masses. With the rise of non-invasive diagnostic tools like transient elastography (FibroScan) and specific blood panels, the overall incidence of liver biopsy has decreased, but it remains the "gold standard" for staging certain liver conditions.

\section*{Aetiology and Pathophysiology}
A liver biopsy is indicated to investigate various pathological processes affecting the liver:

\begin{itemize}
    \item \textbf{Fibrosis and cirrhosis staging:} Quantifying the amount of scar tissue resulting from chronic inflammation.
    \item \textbf{Inflammatory activity grading:} Determining the degree of active necroinflammation in conditions like chronic hepatitis.
    \item \textbf{Infiltration and storage diseases:} Diagnosing lipid deposition (steatosis), iron overload (haemochromatosis), or copper accumulation (Wilson's disease).
    \item \textbf{Unexplained abnormalities:} Investigating persistent elevations in liver enzymes when non-invasive tests are inconclusive.
\end{itemize}

\section*{Clinical Features}
Patients preparing for or recovering from a liver biopsy typically present with features related to their underlying liver condition or the procedure itself:

\begin{itemize}
    \item \textbf{Underlying disease symptoms:} Jaundice (yellowing of the skin and eyes), fatigue, abdominal swelling (ascites), or easy bruising.
    \item \textbf{Procedural anxiety:} Common concern regarding pain, needle insertion, or potential complications.
    \item \textbf{Post-procedure pain:} A mild, dull ache in the right upper abdomen or right shoulder (referred pain via the phrenic nerve) lasting a few hours.
    \item \textbf{Haemodynamic stability:} Close monitoring of pulse and blood pressure post-procedure to detect internal bleeding.
\end{itemize}

\section*{Differential Diagnosis}
Before performing a liver biopsy, other potential causes of hepatic signs and symptoms must be distinguished and evaluated:

\begin{itemize}
    \item \textbf{Biliary tract obstruction:} Gallstones or biliary strictures causing abnormal liver function tests.
    \item \textbf{Congestive hepatopathy:} Liver congestion due to right-sided heart failure.
    \item \textbf{Drug-induced liver injury (DILI):} Liver inflammation caused by medications, herbal supplements, or toxins.
    \item \textbf{Non-alcoholic fatty liver disease (NAFLD):} Investigated non-invasively before considering tissue biopsy.
\end{itemize}

\section*{When to Seek Medical Care}
Patients should discuss the necessity and risks of a liver biopsy with their gastroenterologist or hepatologist. After the procedure, patients must seek immediate emergency medical care if they experience signs of bleeding or infection.

\textbf{Contact your local emergency services for:}
\begin{itemize}
    \item Severe, worsening pain at the biopsy site, in the abdomen, or in the right shoulder
    \item Shortness of breath, rapid breathing, or chest pain
    \item Dizziness, lightheadedness, fainting, or severe weakness (signs of internal bleeding)
    \item Significant bleeding or fluid leaking from the biopsy site
    \item High fever, chills, or persistent shivering
\end{itemize}

\section*{Diagnosis and Investigations}
Before a liver biopsy is performed, several safety investigations are mandatory to minimise risks:

\begin{itemize}
    \item \textbf{Coagulation profile:} Checking international normalised ratio (INR) and prothrombin time (PT) to ensure the blood clots normally.
    \item \textbf{Full blood count (FBC):} Assessing platelet count to verify there is no severe thrombocytopenia.
    \item \textbf{Imaging studies:} Ultrasound or CT scan to map the liver anatomy, locate masses, and plan the needle pathway.
    \item \textbf{Histopathological examination:} Microscopic analysis of the retrieved liver core by a pathologist using special stains.
\end{itemize}

\section*{Management}
The management surrounding a liver biopsy focuses on preparation, procedure execution, and post-operative monitoring:

\begin{itemize}
    \item \textbf{Medication review:} Stopping blood-thinning medicines (e.g., warfarin, aspirin, clopidogrel, or NSAIDs) for 5 to 7 days before the procedure under medical guidance.
    \item \textbf{Fasting:} Requiring the patient to have nothing by mouth (fast) for 6 to 8 hours before the biopsy.
    \item \textbf{Procedural support:} Using local anaesthetic at the insertion site and sometimes conscious sedation to keep the patient comfortable.
    \item \textbf{Post-procedure observation:} Requiring the patient to lie on their right side for 1 to 2 hours, followed by flat bed rest and hourly vital sign checks for a total of 4 to 6 hours.
    \item \textbf{Restricted activity:} Avoiding heavy lifting, vigorous exercise, or straining for 1 to 2 weeks after the procedure.
\end{itemize}

\section*{Complications}
\begin{itemize}
    \item Internal bleeding (haemorrhage), the most serious complication, which may require blood transfusion or surgical intervention
    \item Significant pain requiring analgesia, often referred to the right shoulder
    \item Accidental puncture of other organs, such as the lung (pneumothorax), gallbladder, or kidney
    \item Infection (peritonitis or localized abscess) at the insertion site or within the abdominal cavity
    \item Bile leak leading to biliary peritonitis
\end{itemize}

\section*{Prognosis and Outcomes}
The prognosis for recovery from the liver biopsy procedure itself is excellent, with most patients resuming normal light activities within a few days. The clinical outcome depends entirely on the disease diagnosed by the biopsy. The biopsy results provide critical staging and diagnostic information, allowing clinicians to tailor specific, life-saving treatment plans for conditions like autoimmune hepatitis, chronic hepatitis B or C, or fatty liver disease.

\section*{Prevention and Self-Care}
\begin{itemize}
    \item Follow all pre-operative instructions regarding fasting and medication suspension carefully
    \item Keep the biopsy dressing clean and dry for at least 24 hours
    \item Avoid strenuous exercise, sports, and heavy lifting for the recommended recovery period
    \item Arrange for a responsible adult to drive you home and stay with you overnight after the procedure
    \item Monitor the puncture site daily for signs of infection, such as redness, swelling, or discharge
\end{itemize}

\section*{Sources and Further Reading}
The following organisations provide reliable information on liver biopsy:

\begin{itemize}
    \item Gastroenterological Society of Australia. \textit{Patient guide to liver biopsy procedures and preparation.} https://www.gesa.org.au/
    \item Mayo Clinic. \textit{Clinical overview of liver biopsy types, risks, and recovery.} https://www.mayoclinic.org/
    \item NHS. \textit{Information on what happens during and after a liver biopsy.} https://www.nhs.uk/
    \item Healthdirect Australia. \textit{Liver biopsy overview and patient guidance.} https://www.healthdirect.gov.au/
    \item MSD Manual. \textit{Professional reference on liver biopsy techniques and findings.} https://www.msdmanuals.com/
\end{itemize}
